\scanpage{330}
\mspage{69}
% archive sheet number: 329
% ms. 1–5
$\beta(c)$, pour une $\Sigma$-flèche de type $\tau$, \inserted{de \unsure{long.|longueur} $n$,} en inversant toutes les flèches rétrogrades (\deleted{\unsure{de}}ce qui est possible, puisqu'elles sont dans $\Sigma$), de façon à obtenir un chemin \inserted{$c^{\unsure{\leftrightarrows|\leftrightsquigarrow}}$} de type $\Delta_+^n$, et en prenant le composé (dans $X$) des flèches de ce chemin $c^{\leftrightarrows}$.
% ms. 5–8
Il est évident donc que $\beta\alpha=\mathrm{id}$, il faut prouver que l'on a aussi $\alpha\beta=\mathrm{id}$, i.e. que dans $\Pi_\Sigma$, les flèches $u=\int c^{\leftrightarrows}$ et $c$ sont égales, i.e. $u$ et $c$ sont homotopes dans $\underline{\mathrm{Ch}}_\Sigma(a,b)$.
% ms. 9
Considérons le chemin $c^{\leftrightarrows}$, de type $\Delta_+^n$ :
% ms. 10
\diagram{\begin{tikzcd}[ampersand replacement=\&,column sep=small]
a=a_0\arrow[r,"u_1"] \& a_1\arrow[r,"u_2"] \& a_2 \& a_{n-1}\arrow[r,"u_n"] \& a_n=b
\end{tikzcd}}
% ms. 11–14
et soit $u_i$ ($i\in[1,n]$) une flèche de ce chemin, correspondant à une flèche de transition \inserted{$\underline{u}_i$} \uline{rétrograde} de $\tau$, donc par construction
% ms. 14
\[u_i=v_{i\deleted{\unsure{u}}}^{-1}\qquad\text{dans } X,\quad\text{avec }v_i\in\Sigma\]
% ms. 15–19
la flèche d'indice \deleted{\unsure{un}} $i$ dans le chemin initial $\tau$ de \deleted{\unsure{$f$}} type $\tau$. \deleted{Or on prouve que \unsure{alors}} \deleted{et cela implique que \unsure{dans les}} Si $\underline{u}_i$ par contre est \deleted{$\mathrm{ch}(u_i)=(\mathrm{ch}(v_i))^{-1}\text{ dans }\Pi_\Sigma(X)$} directe, on a $u_i=v_i$. \deleted{\unsure{Or on a}} (NB On désigne par $v_i$ les $\unsure{i^e|i}$ flèches dans le chemin $c$.) Or on a
% ms. 20
\[\begin{aligned}
&\deleted{\unsure{c(u)=c(u_i)}}\\
&c=c_n*c_{n-1}\cdots c_1\qquad\text{dans }\Pi_\Sigma(X)
\end{aligned}\tag*{$(+)$}\]
% ms. 21–26
et les facteurs $c_i$ ont soit \deleted{donnés par}\inserted{la classe d'homot. de} $\mathrm{ch}(\unsure{v_i|u_i})$ \inserted{$=\mathrm{ch}(u_i)$} (cas où la flèche de transition $\underline{u}_i$ de $\tau$ est directe), soit celle de \deleted{$\mathrm{ch}(u_i)^{\unsure{\circ|o}}$,} du chemin \uline{opposé} au chemin $\mathrm{ch}(v_i):a_i\xrightarrow{v_i}a_{i-1}$ de type $\Delta_+$, donc le chemin $\mathrm{ch}(v_i)^\circ:a_{i-1}\xleftarrow{v_i}a_i$ de type $\Delta_-$. Je dis que
% ms. 27
\[\tikz[baseline=(rel.base)]{\node[inner sep=2pt] (rel) {$\mathrm{ch}(v_i)^\circ\deleted{\unsure{=\mathrm{ch}(v_i)^{-1}}}$};\draw (rel.south west) -- (rel.south east);\draw (rel.south west) -- ++(0,.18);\draw (rel.north west) -- ++(.48,0);}\inserted{=\mathrm{ch}(u_i)}\qquad\text{\uline{dans} }\Pi(X)\tag*{(1)}\]
% ms. 28–30
ce qui prouve que dans $\Pi_\Sigma(X)$, \deleted{le produit} $c$ est égal au composé des $\mathrm{ch}(u_i)$ $(i\in[1,n])$ donc à $c^{\leftrightarrows}$, \deleted{et \unsure{cela} $c^{\leftrightarrows}$ \unsure{est égal à l'objet}} \inserted{et \unsure{il faut|il nous}} prouver de plus que
% ms. 30
\[c^{\leftrightarrows}=\mathrm{ch}(u)\qquad (u=u_nu_{n-1}\cdots u_{\unsure{1|p}})\quad\text{dans }\Pi(X).\tag*{(2)}\]
